\chapter{Introduction}
\label{ch:essay-intro}

A spoken word, a gesture, and a musical phrase share no common physical substrate. The acoustic pressure wave that carries speech, the motor trajectory that produces a pointing motion, and the harmonic structure that constitutes a melodic contour differ not merely in details of their realization but in the kind of physical system through which each is instantiated. The puzzle arises when one observes that these heterogeneous signals may nonetheless participate in the communication of related or equivalent meanings. This becomes theoretically puzzling only if the observable signal is assumed to be identical with the structured object being communicated. Once that assumption is rejected, the problem shifts from explaining how different physical substrates could mean the same thing to determining which structural relations survive the projection from latent structure into observable signal, which are altered by the constraints of the physical channel, and which are destroyed so completely that reconstruction becomes impossible.

The thesis developed here is that the observable is a potentially lossy projection of latent structure. The question is therefore not whether reconstruction is possible in general — clearly some projections preserve more structure than others, and some preservation is sufficient for successful communication — but rather what survives projection and what does not. Four theorems developed elsewhere in this research program, proven in the monograph \emph{Admissible Reconstruction}, establish the formal apparatus needed to answer this question precisely. Each theorem addresses a different aspect of the same underlying admissibility relation, examined at different points in its behavior.

The first question is when two representations, possibly realized through entirely different physical or computational substrates, ought to be considered the same object. Representation Invariance, proven in Chapter~1 of \emph{Admissible Reconstruction}, establishes the conditions under which physically different representations count as semantically equivalent. The answer turns on whether the admissibility structures induced by the two representations are related by an appropriate isomorphism, independent of the substrate through which either is realized. The theorem therefore provides a criterion for deciding when substrate differences are semantically irrelevant.

The second question is what it means for a trajectory through admissibility structure to remain coherent. Admissibility Preservation, proven in Chapter~2, addresses this directly. Coherence is not an additional property imposed on top of admissibility; it is what nontriviality of the admissibility relation amounts to when examined along a trajectory rather than at a single state. A trajectory degenerates, and thereby fails to be coherent, exactly when the reachable set collapses in a specific technical sense defined by the theorem. The question is therefore whether the relevant admissibility relation continues to discriminate admissible from inadmissible continuations at each step.

The third question concerns the conditions under which a scalar measure of inadmissibility decreases along a trajectory. Monotonic Relaxation, proven in Chapter~3, establishes that such decrease occurs along trajectories following repair dynamics specifically — transitions selected to increase the volume of the reachable set — but not along admissible trajectories in general. The theorem thereby identifies the precise sense in which certain restorative processes exhibit convergence, while carefully distinguishing this claim from a separate ordering relation developed elsewhere in this program.

The fourth question is when it becomes impossible in principle, not merely difficult in practice, to recover which of two histories actually occurred from the evidence either would have left behind. Non-Identifiability, proven in Chapter~4, establishes both an elementary direction — that when two distinct histories produce identical observables, no single-valued reconstruction function can return the correct distinct history in both cases — and a structural direction, showing that such collisions are not pathological exceptions but a generic feature of observation maps constrained to preserve only reachable-set structure. The theorem additionally provides a diagnostic criterion for distinguishing cases where reconstruction failure reflects a correctable observational error from cases where it reflects a fundamental structural limit.

These four results provide the formal foundation for the cross-modal investigation that follows. Representation, coherence, repair, and the limits of recovery are not four separate phenomena requiring four independent theories. They are four consequences of a single admissibility relation, examined respectively at a single state, along a trajectory, under restorative dynamics, and under lossy projection.
